\VerdierPhysicalPageAnchor{0021}
the morphism of complexes of sheaves:

\VerdierSemanticAnchor{intro-1-flat-quasi-isomorphism}
\[
F\otimes\rho:
F\otimes_{\mathbb Z}L_1^{\bullet}
\longrightarrow
F\otimes_{\mathbb Z}L_2^{\bullet}
\qquad\text{,}
\]
is a \emph{quasi-isomorphism}.

\VerdierSemanticAnchor{intro-1-coherent-tensor-system}
Thus one obtains a coherent system of complexes of sheaves
$F\otimes_{\mathbb Z}L^{\bullet}$ and quasi-isomorphisms which depends
functorially on the sheaf $G$, in a sense that can be made precise. In any
case, the cohomology objects of the complexes
$F\otimes_{\mathbb Z}L^{\bullet}$ ($L^{\bullet}$ a flat resolution of $G$) are
defined up to unique isomorphism and depend functorially on $G$. The resulting
functors have the desired properties and deserve the name left derived
functors of the functor ``tensor product with $F$''.

\VerdierSemanticAnchor{intro-1-tor-sheaves}
Of course, to calculate the torsion sheaves
$\operatorname{Tor}^{\mathbb Z}(F,G)$, one can take the sheaf associated with
the presheaf
$U\longmapsto\operatorname{Tor}^{\mathbb Z}(F(U),G(U))$ ($U$ open in $E$).
But the method we use provides a complex of sheaves well defined up to
quasi-isomorphism whose cohomology objects are
$\operatorname{Tor}^{\mathbb Z}(F,G)$ and $F\otimes_{\mathbb Z}G$, additional
information essential for handling these functors conveniently, as we shall
see below. This method extends immediately to ringed topoi and makes it
possible in this setting to define the left derived functors of the functor
``tensor product of sheaves''.

\VerdierSemanticAnchor{intro-2}
\noindent\textbf{2.}\quad
In the study of the derived functors of a composite of two functors, of
associativity properties, and of relations of Künneth type, one is led to
extend the formalism of derived functors to the case where the argument is no
longer merely an object of the category under study, but a complex of objects
of that category. This extension was carried out in [1]. Let us recall its
main outlines. Let $F:\mathcal A\to\mathcal B$ be an additive functor between
two abelian categories such that $\mathcal A$ has enough injective objects.
Let $Y^{\bullet}$ be a complex of objects of $\mathcal A$. In \emph{loc. cit.}
one defines the ``injective resolutions'' of the complex $Y^{\bullet}$, which
we call \emph{Cartan--Eilenberg injective resolutions} of $Y^{\bullet}$.

\VerdierSemanticAnchor{intro-2-ce-resolution}
Such a resolution is, by definition, a double complex
$I^{\bullet\bullet}$ whose components are injective objects of $\mathcal A$,
equipped with an augmentation

\VerdierSemanticAnchor{intro-2-ce-augmentation}
\[
Y^{\bullet}\xrightarrow{\varepsilon}I^{\bullet\bullet}
\]
that induces injective resolutions (in the usual sense) of the components,
the boundaries, and the cohomology objects of the complex $Y^{\bullet}$. These
resolutions are unique up to homotopy (of double complexes). One then forms
the double complex $F(I^{\bullet\bullet})$, and then the associated simple
complex $\int F(I^{\bullet\bullet})$, which likewise is determined
